\documentclass[a4paper,12pt,oneside]{article}

\usepackage[T2A]{fontenc}
\usepackage[utf8]{inputenc}
\usepackage[russian]{babel}

\usepackage{lmodern}
\usepackage{microtype}

\usepackage{amsmath,amssymb,amsthm,mathtools,mathrsfs}
\usepackage{dsfont}

\usepackage{geometry}
\geometry{top=25mm,left=25mm,right=20mm,bottom=25mm}

\usepackage{graphicx}

\usepackage{enumitem}
\setlist[enumerate]{leftmargin=*, itemsep=0.35\baselineskip, topsep=0.35\baselineskip}
\setlist[itemize]{leftmargin=*, itemsep=0.2\baselineskip, topsep=0.25\baselineskip}

\newlist{alphlist}{enumerate}{1}
\setlist[alphlist]{label=\textbf{\alph*)}, leftmargin=*, align=left}

\newlist{numlist}{enumerate}{1}
\setlist[numlist]{label=\textbf{\arabic*)}, leftmargin=*, align=left}

\newcommand{\Task}[1]{%
  \subsection*{#1}%
  \vspace{-0.2\baselineskip}%
}

\setlength{\parindent}{1.25em}
\setlength{\parskip}{0pt}

\begin{document}

\begin{center}
    {\Large\bfseries Множества и операции над ними}\\[1mm]
\end{center}
\vspace{2mm}
\hrule
\vspace{4mm}

\Task{Задание 1}
Задайте формулой:
\begin{numlist}
    \item $\{0, 1, 1, 2, 3, 5, 8, \dots\}$
    \item $\{19, 38, 57, 76, 95, \dots\}$
    \item $\{2, 6, 18, 54, 162, \dots\}$
\end{numlist}

\Task{Задание 2}
Известно, что $y \in A \cup B$. Верно ли, что $y \in A$?

\Task{Задание 3}
В стране $N$ парламентские фракции образованы так, что для любых двух фракций $A$ и $B$ (необязательно различных) верно, что $\overline{A \cup B}$ является фракцией. Докажите, что $\forall A,B \Rightarrow A \cup B$ — тоже фракция.

\Task{Задание 4}
Изобразите на графиках следующие множества точек:
\begin{numlist}
    \item $\{1,2,3\} \times (1; 3]$
    \item $[1; 5) \times (1; 4] \setminus \{(2, 3)\}$
    \item $[1; 7] \times (1;5] \setminus (1;4] \times (1;3)$
\end{numlist}

\Task{Задание 5}
Воспользовавшись свойствами операций над множествами, выяснить, верны ли тождества для любых $A, B, C$:
\begin{numlist}
    \item $A \setminus B = A \,\Delta\, (A \cap B)$
    \item $A \cap (B \setminus C) = (A \cap B)\setminus (A \cap C)$
    \item $\overline{A} \cup (A \cap C) = \overline{A} \cup C$
    \item $A \setminus (B \setminus C) \neq (A \setminus B) \setminus C$
    \item $A \cap (B \cup C) = (A \cap B) \cup (A \cap C)$
    \item $A \cup (B \Delta C) = (A \cup B) \Delta (A \cup C)$
    \item $(A \Delta B) \Delta C = A \Delta (B \Delta C)$
    \item $A \cup (A \cap B) = A$
    \item $(A \cap B) \cup (A \cap C)$
\end{numlist}

\Task{Задание 6}
Задайте множество аналитически по диаграммам Эйлера–Венна:

\includegraphics[width=0.35\linewidth]{image1.png}
\includegraphics[width=0.35\linewidth]{image2.png}
\includegraphics[width=0.35\linewidth]{image3.png}

Упростите полученные формулы.

\Task{Задание 7}
Найдите мощность булеана у следующих множеств:
\begin{numlist}
    \item $\{1, 2, 3\}$
    \item $\{1, a, 2, b\}$
    \item $\{a, \{a\}, \{a, b\}, b\}$
    \item $\{x \mid P(x) = 0,\; P(x) = x^2 + 21x + 550,\; x \in \mathbb{C}\}$
\end{numlist}

\end{document}
